\chapter{Matching designs and the concrete lower bounds}
\label{tx:app:matching}

The examples use a lower bound for a concrete rectangular system.  Its
proof illustrates why a functional excluding a Nullstellensatz certificate
does not by itself exclude a Polynomial Calculus derivation.  To handle
degree falls we construct a whole space of functionals, together with an
extension operator between successive degrees.
The use of designs as linear obstructions to Nullstellensatz proofs
has a direct antecedent in Buss~\cite{BussDesign}. The argument here
also proves the extension property needed to handle ordinary-PC lines.
Here the target is $1$. The capital $D$ below is the proposed proof-degree
budget, rather than the input-degree bound of a general template.
The integers $N,M$ describe row and column counts in the auxiliary boards.

\section{Why this separate example proof is informative}
The general linear theorem used an exact operation on whole consequence
spaces. The concrete rectangular family gives a second way to see why
ordinary-PC degree falls must be handled honestly. Here we will use
linear functionals instead of a group-algebra projector.

A single functional annihilating the permitted axiom products excludes
an NS certificate. It need not remain valid when a cancelled low-degree
polynomial is multiplied in PC. The remedy is not to announce that NS
and PC are the same. We construct a whole space of designs and prove
that every lower-degree design extends. The PC invariant will quantify
over \emph{all} designs. This universal quantifier is what survives a
multiplication after cancellation.

The proof has four steps: choose a terminating signed reduction of
matchings, prove independence of its integral coordinates in a square,
compress complete histories to obtain a binomial identity in rectangles,
and use that identity to extend designs. The last step applies the
functionals to every legal PC line. We give the complete construction
in this chapter, rather than referring to a pigeonhole lower-bound paper.

For positive integers $r\leq m$, let $B_{r,m}$ consist of the equations
\[
 \sum_{j=1}^{m}x_{ij}-1=0\quad(i\in[r]),\qquad
 \sum_{i=1}^{r}x_{ij}-1=0\quad(j\in[m]),
\]
the Boolean equations $x_{ij}^{2}-x_{ij}=0$, and the row and column
collision equations
\[
 x_{ij}x_{i\ell}=0\quad(j\ne\ell),\qquad
 x_{ij}x_{hj}=0\quad(i\ne h).
\]
Here $[r]=\{1,\ldots,r\}$.  A Boolean solution would describe a bijection
between the two sides, so a proper rectangle has no solution.  Nevertheless,
in positive characteristic the two linear totality sums can agree because
$m-r$ vanishes in the field.  The following statement identifies a range
in which this obstruction persists through ordinary PC.

\begin{lemma}[Rectangle obstruction]\label{tx:lem:rectangle-obstruction}
Let $K$ be a field of prime characteristic $p$, and let $a,D$ be
nonnegative integers.  If
\[
 1\leq r\leq m,\qquad 2D\leq r,\qquad D<p^a,
 \qquad p^a\mid m-r,
\]
then $B_{r,m}$ has no ordinary-PC refutation of degree at most $D$.
Consequently it has no Nullstellensatz refutation of total-product degree
at most $D$.
\end{lemma}

We give the full construction over the integers before passing to $K$.
An \emph{edge set} is a subset of $[N]\times[M]$.  It is a matching if
no two distinct edges share a row or a column.  Write $\operatorname{rows}P$
and $\operatorname{cols}P$ for its two projections.  A degree-$d$ design
may take values in a coefficient ring: it is a function $\psi$ on edge
sets, zero on nonmatchings, for which
\begin{equation}\label{tx:eq:matching-design}
 \sum_{j=1}^{M}\psi(P\cup\{(i,j)\})=\psi(P),\qquad
 \sum_{i=1}^{N}\psi(P\cup\{(i,j)\})=\psi(P)
 \quad(|P|<d).
\end{equation}
The indices $i\in[N]$ and $j\in[M]$ are arbitrary. For $K$-valued
designs, denote the resulting vector space by $E_d(N,M)$. The signed
coordinate reductions below are constructed integrally and then
interpreted in $K$.  Normalisation, $\psi(\varnothing)=1$, is a separate
condition.  We will construct a normalised member of $E_D(r,m)$ and prove
that, for $D\ge1$, \emph{every} member of $E_{D-1}(r,m)$ extends to $E_D(r,m)$ while
preserving its values through degree $D-1$.

To interpret the table, associate an edge set $P$ with the squarefree
monomial $x_P=\prod_{e\in P}x_e$. Its proposed value is $\psi(P)$.
For an actual matching assignment this value would be its containment
indicator; a design instead imposes only the stated linear consistency
conditions in the tested degree window. A normalised design can be read
as a \emph{pseudoexpectation} in the following algebraic sense: a linear
assignment of values to monomials satisfying those tests and giving the
constant monomial value one, without a probability distribution or a
positivity assumption. Its extension to ordinary monomials is spelled
out in the last section. Unnormalised designs are also needed for the
extension argument.

\section{A degree-one design which is not a solution}
Over $\mathbb F_2$, consider four rows and six columns. Give the six
columns respectively the selected row labels $1,1,1,2,3,4$, and let
$\psi(\{(i,j)\})$ be one on those six selected edges and zero elsewhere.
Let $\psi(\varnothing)=1$. Every column sum is one. The four row
sums are $3,1,1,1$, all one in the field. Thus these values satisfy the
degree-one totality conditions, although the selected edges do not form
an actual bijection: three columns use the same row.

This is an honest linear witness, not a Boolean model. Degree one has
not yet tested the quadratic collision axioms. Higher-degree designs
must also assign compatible values to matching edge sets. The extension
construction below supplies these values exactly in the range of the
rectangle obstruction, without division by a factorial.

\begin{checkpoint}
An actual solution evaluates every polynomial multiplicatively. A design
is a consistent linear response only to the tested degree window. It
can exist for an inconsistent rectangle and still rule out every proof
within that window.
\end{checkpoint}

\section{A terminating signed reduction}
The reduction will replace one matching by a formal integer linear
combination of others. The word \emph{signed} allows negative
coefficients; it does not count matchings with positive weights.
The required identity is that every design gives the same value before
and after a replacement. Keeping integer coefficients first lets us
use the same identities after reducing modulo any characteristic.


A chain in a matching is a subset whose row coordinates and column
coordinates both increase.  Let $h(P)$ be the maximum chain length, and
write $P_{\leq b}$ for the edges in rows at most $b$.  A matching $P$ is
\emph{terminal} if, for every positive integer $b$,
\[
 h(P_{\leq b})\leq b-|\operatorname{rows}(P_{\leq b})|,
\]
and the corresponding inequality holds for column prefixes.  The right
side counts unused labels in the prefix.  This definition depends only
on the labelled matching, not on the ambient rectangle.  Prefixes beyond
its last occupied row or column cannot introduce a new failure: their
left side stays fixed and their right side increases.

For a nonterminal matching, use the least failed row prefix if one exists;
otherwise use the least failed column prefix.  In the selected prefix
choose a longest chain, breaking ties by the lexicographically greatest
list of its edges, with the edges themselves in row-first lexicographic
order.  In a column step use this rule after transposition.  Thus all
choices are deterministic.

Consider a row step with failed prefix $b$ and chosen chain $A$.  Put
\[
 P_0=P\setminus A,\qquad S=\operatorname{cols}A,\qquad
 B=\{i\leq b:i\notin\operatorname{rows}P\},
\]
\[
 U=\operatorname{rows}A\ \cup\
       \{b<i\leq N:i\notin\operatorname{rows}P\}.
\]
The sets $B$ and $U$ are disjoint and partition all rows unused by $P_0$.
Moreover, $|B|<|A|$, because the selected prefix fails.  A longest chain
in the first failed prefix ends in row $b$.  Indeed, if it ended earlier,
the same chain would belong to the preceding prefix, whereas the number
of unused labels cannot decrease when one adds the next row.  That
preceding prefix would already fail.

For sets of columns $S$ and rows $V$, a full injection from $S$ into $V$
means a matching $J$ with $\operatorname{cols}J=S$ and
$\operatorname{rows}J\subseteq V$.  A partial injection has only
$\operatorname{cols}J\subseteq S$.  Replace $P$ by
\[
 \sum_{\substack{J\text{ partial injection}\\S\to B}}
       (-1)^{|J|}[P_0\cup J]
 \ -\!
 \sum_{\substack{J\text{ full injection}\\S\to U\\J\ne A}}
       [P_0\cup J],
\]
omitting colliding terms.  Square brackets here denote formal basis
elements, not polynomial equivalence classes.  The transposed rule defines
a column step.

To verify the identity for a degree-$d$ design with $|P|\leq d$, repeatedly
apply its column equation to the columns of $S$.  At each application
the partial matching has fewer than $|P|$ edges, so the degree restriction
is satisfied.  An occupied row contributes zero by collision vanishing.
For a still unassigned column the sum over available rows in $U$ is
therefore the previous value minus the sum over rows in $B$.  Expanding
these successive differences gives
\[
 \sum_{J:S\to U}\psi(P_0\cup J)
  =\sum_{A'\subseteq S}(-1)^{|A'|}
       \sum_{J:A'\to B}\psi(P_0\cup J).
\]
Every surviving term is an injection, because collisions have value zero.
The right side contains each partial injection exactly once, indexed by
its column set.  Isolating the term $J=A$ on the left proves the stated
reduction.  The column step follows by interchanging rows and columns in
the design equations.  No assertion that the deterministic reduction
itself commutes with transposition is needed.

We next prove termination, including the equal-degree case.  For the
fixed grid $G=[N]\times[M]$, set $H=N+M$, $c=NM+1$, and define
\[
 w(i,j)=c^{H-i-j},\qquad W(P)=\sum_{(i,j)\in P}w(i,j),\qquad
 \rho(P)=|P|(W(G)+1)+W(P).
\]
Every lower term has fewer edges, since an injection into $B$ has at most
$|B|<|A|$ edges.  Thus its $\rho$ is smaller irrespective of its weight.
For a full injection $J\ne A$ into $U$, let $e$ be the first edge of $A$
not belonging to $J$, in increasing row order.  Every earlier chain edge
still belongs to $J$.  If $x\in J\setminus A$ had row smaller than that
of $e$, its row would be an earlier chain row already occupied by its
retained edge; this contradicts the matching condition.  A row from the
other part of $U$ is greater than $b$ and hence cannot be smaller.
Likewise, a column of $x$ smaller than the column of $e$ belongs to an
earlier retained chain edge, again giving a collision.  Consequently
\[
 e_1\leq x_1,\qquad e_2\leq x_2,\qquad e_1+e_2<x_1+x_2.
\]
At least one inequality is strict because $x\ne e$.  Each newly added
edge has weight at most $w(e)/c$, and there are at most $NM<c$ such
edges.  Their total weight is strictly smaller than $w(e)$.  Cancelling
the common edges of $A$ and $J$ shows $W(P_0\cup J)<W(P)$.
The identical argument applies to a column step, since $i+j$ is
unchanged by transposition.  Every child therefore has strictly smaller
nonnegative integer rank $\rho$.  The tree is finitely branching, and
every path has length at most its initial rank.  Hence the complete
signed reduction tree is finite.

For terminal $T$, let $C_{N,M}(P,T)\in\mathbb Z$ be the sum, over paths
from $P$ to $T$, of the products of their step signs; set it to zero for
colliding $P$.  The empty path gives
\[
 C_{N,M}(S,T)=\delta_{S,T}\quad(S\text{ terminal}),\qquad
 C_{N,M}(P,T)=0\quad(|T|>|P|).
\]
Induction on $\rho(P)$ using the reduction identity gives, for every
degree-$d$ design,
\begin{equation}\label{tx:eq:terminal-expansion}
 \psi(P)=\sum_{\substack{T\text{ terminal}\\|T|\leq d}}
              C_{N,M}(P,T)\psi(T)\qquad(|P|\leq d).
\end{equation}
Thus the terminal values determine every design in the relevant range.
It remains to show that these coordinates themselves satisfy the required
totality equations and can be restricted from a square to a rectangle.

For a small illustration of the preceding reduction, in a $2$-by-$2$
square the matching $\{(1,2)\}$ fails the first row prefix. Its chosen
chain is that one edge, $B$ is empty, and $U=\{1,2\}$. The rule gives
\[
 [\{(1,2)\}]\longmapsto[\varnothing]-[\{(2,2)\}].
\]
Here $\{(2,2)\}$ is terminal. For a design of degree at least one,
the corresponding identity is simply the column sum $\psi(\{(1,2)\})+\psi(\{(2,2)\})=\psi(\varnothing)$.
Thus a terminal coordinate can be negative; it is not a probability.

\section{Integral independence in a square}
An indicator here is a function on actual permutations of the square:
it takes value one when the permutation contains the given matching,
and zero otherwise. The argument is \emph{integral} because it uses
integer coefficients throughout. Here it proves independence after
passing to any coefficient ring in the theorem, as well as over
the integers. The signed moments below isolate a coefficient
with multiplier $1$ or $-1$, both invertible in every ring. This is the
reason for proving a unit diagonal rather than a nonzero factorial.


Let $T\subseteq[N]\times[N]$ be terminal.  For an edge $e\in T$, let
its height be the maximum chain length in the southwest downset
$\{f\in T:f_1\leq e_1,\ f_2\leq e_2\}$.  A longest such chain ends
at $e$: otherwise $e$ can be appended, because distinct edges of a
matching have distinct row and column coordinates.  If $e$ is southwest
of $f$ with both coordinates strictly smaller, appending $f$ raises the
height.  It follows that every equal-height level $T_t$ is decreasing:
as its row coordinate increases, its column coordinate decreases.
The levels partition $T$.

Let $h=h(T)$.  The terminal inequality at prefix $N$ implies that there
are at least $h$ unused rows and $h$ unused columns.  Let $a_t$ and $b_t$
be the $t$th smallest unused row and column, respectively, for
$1\leq t\leq h$.  These anchors precede every coordinate in $T_t$.
For example, if $e\in T_t$, a chain of length $t$ ending at $e$ belongs
to its row prefix.  Terminality requires at least $t$ unused rows in
that prefix; row $e_1$ is occupied, so the $t$th unused row is strictly
less than $e_1$.  The column argument is identical.  Form the paired blocks
\[
 A_t=\operatorname{rows}T_t\cup\{a_t\},\qquad
 B_t=\operatorname{cols}T_t\cup\{b_t\}.
\]
They are disjoint on each side, their paired sizes agree, and their
minima are $a_t,b_t$.  Their complements also have equal sizes.  Fix the
increasing bijection on the complements, and an increasing base bijection
$A_t\to B_t$ in each block.  Vary each base bijection by a permutation
$\sigma_t$ of $A_t$.  The union of all these bijections is an actual
permutation $\pi_\sigma$ of $[N]$.

For any edge set $P$, define the integer signed moment
\[
 M_T(P)=\sum_{\sigma}
     \left(\prod_{t=1}^{h}\operatorname{sgn}\sigma_t\right)
     \boldsymbol{1}_{P\subseteq\pi_\sigma}.
\]
If a block has two rows absent from $\operatorname{rows}P$, interchange
their images.  This is a fixed-point-free involution on the choices,
changes the sign, and preserves the containment condition.  Its terms
cancel over $\mathbb Z$.  In particular, $M_T(P)=0$ when $|P|<|T|$:
otherwise every block would use at least $|A_t|-1=|T_t|$ rows, forcing
$|P|\geq\sum_t|T_t|=|T|$.

The diagonal is a unit.  In each block the edges of $T_t$ prescribe the
bijection on every row except $a_t$; its remaining value is necessarily
$b_t$.  Thus exactly one choice $\sigma$ contains $T$, and
\[
 M_T(T)\in\{1,-1\}.
\]
This is a single signed contribution; no factorial or positive integer
multiplicity must be inverted.

For equal-degree triangularity, suppose $|P|=|T|$ and $M_T(P)\ne0$.
There is a choice containing $P$, so $P$ is a matching.  Cancellation
shows that at most one row in each block is unused.  Counting over the
disjoint blocks forces exactly $|T_t|$ edges in block $A_t\times B_t$,
one unused row in every block, and no edge in the complement.  Write
$P_t$ for this block matching.

Among all matchings of size $|T_t|$ in $A_t\times B_t$, the decreasing
matching $T_t$ has uniquely least weight $W$.  Here is a direct proof.
Choose a weight minimiser $Q$ from this finite set.  If it uses the
minimum row $a_t$, move that edge to the unique unused row.  This is a
legal matching, has the same cardinality, and strictly smaller weight,
because its row coordinate has increased.  Thus $a_t$ is unused.
The same argument shows that the minimum column $b_t$ is unused.
If $Q$ contains an increasing pair $(i,j),(i',j')$ with
$i<i'$ and $j<j'$, replace it by the crossed pair
$(i,j'),(i',j)$.  Both added edges have coordinate sum strictly larger
than that of $(i,j)$, and their combined weight is smaller than its
weight, since $c=N^2+1>2$ when such a pair exists.  The unchanged context
is disjoint from the crossed pair.  This again contradicts minimality.
Therefore $Q$ is decreasing, uses precisely $A_t\setminus\{a_t\}$
and $B_t\setminus\{b_t\}$, and must equal $T_t$: on two fixed ordered
sets there is exactly one order-reversing bijection.  Any distinct
matching with the same minimum weight would also be a minimiser and
would give the same conclusion.  Hence
\[
 W(P_t)\geq W(T_t),\quad\text{with equality only if }P_t=T_t.
\]
Summing over the disjoint blocks proves that
\[
 |P|=|T|,\quad P\ne T,\quad M_T(P)\ne0
 \quad\Longrightarrow\quad W(P)>W(T).
\]
Together with the lower-degree cancellation, this says that $M_T(P)=0$
for every $P\ne T$ with $\rho(P)\leq\rho(T)$.

Now suppose a finite linear combination of terminal matching indicators
vanishes at every permutation of $[N]$, with coefficients in any
commutative ring $R$.  Among its nonzero coefficients choose a terminal
$T$ of greatest rank $\rho$.  Apply the integer signed moment $M_T$ to
this relation, then map the integers to $R$.  All other nonzero terms
vanish by triangularity; the result is the coefficient of $T$ multiplied
by $1$ or $-1$.  This contradicts its being nonzero.  The terminal
indicators are therefore independent over every commutative ring,
including fields of characteristic two and rings with zero divisors.

Each square permutation satisfies the design equations in every degree.
Equation~\eqref{tx:eq:terminal-expansion}, applied to its containment
indicator, expresses every matching indicator in terminal coordinates.
For a fixed row $i$, the permutation identity
\[
 \sum_{j=1}^{N}\boldsymbol{1}_{P\cup\{(i,j)\}\subseteq\pi}
       =\boldsymbol{1}_{P\subseteq\pi}
\]
holds because $\pi$ has exactly one edge in that row.  Expand all terms
in terminal coordinates and use the integral independence just proved.
For each terminal $T$ this yields
\[
 \sum_{j=1}^{N}C_{N,N}(P\cup\{(i,j)\},T)=C_{N,N}(P,T).
\]
The column identity follows in the same way from the unique edge in a
column.  Colliding sets have zero coordinates and zero permutation
indicators, so they are included.  Thus every square coordinate is an
unnormalised design, in every degree, over the integers.

\section{Whole-history compression and the marked binomial identity}
The last matching of a path need not remember every row visited by
that path. We must group paths using their \emph{entire} histories.
Once the total signed weight is the same for every chosen $k$-element set
of appended rows, summing over those sets produces $\binom{L}{k}$.
The mark ensures that $k$ is positive, even if its initial edge is
later deleted. This will remove the unwanted appended-row contribution
when the binomial coefficients vanish in $K$.


Fix $P,T\subseteq[r]\times[m]$, with $P$ a matching and $T$ terminal, and a column
$j\notin\operatorname{cols}P$.  Append $L$ rows, and consider all pairs
consisting of a marked initial row $i\in\{r+1,\ldots,r+L\}$ and a signed
reduction path from $P\cup\{(i,j)\}$ to $T$.  The weight of a pair is the
product of its path signs.  Its \emph{history support} is the set of
appended rows appearing anywhere on the path, including the initial
state and edges that are subsequently deleted.

If $|P|+1\leq D$ and $2D\leq r$, every such support is nonempty and has
size at most $|P|+1$.  To prove the latter assertion, give one token to
each initial edge.  In a row rewrite, transport a token by its column
to the unique child edge in that column, or delete it if that column
is absent.  The child uses only columns of the parent, so this describes
every child edge exactly once.  In a column rewrite transport by row,
with the analogous deletion rule.  No deleted token reappears, and no
two surviving tokens represent the same edge.

For any matching $Q$ of at most $D$ edges, failure at row prefix $b$
implies
\[
 b< h(Q_{\leq b})+
       |\operatorname{rows}(Q_{\leq b})|\leq 2|Q|\leq2D\leq r.
\]
Thus every chosen row chain lies in old rows.  An edge already in an
appended row is part of the unchanged context of a row rewrite and is
retained verbatim.  A column rewrite preserves its token's row or
deletes the token.  It follows that a token can encounter at most one
appended row throughout its lifetime.  There are $|P|+1$ initial tokens,
so the full history has at most that many appended rows.  The initial
mark belongs to it even when its token later disappears.

We next verify that histories of the same size carry the same total
integer weight.  Let $S$ be a prescribed $k$-element subset of the $L$
appended rows, and let $f$ be the increasing injection fixing $[r]$ and
sending $r+1,\ldots,r+k$ to the elements of $S$ in order.  For $t\geq1$, set $f(r+k+t)=r+L+t$.
This extends $f$ increasingly and makes the inverse image of the large
row interval $[r+L]$ exactly $[r+k]$.  Rename the row coordinate of every
edge by $f$ and retain its column and step sign.

First, this renaming preserves the deterministic selector.  Every failed
row prefix lies below $r$, as just proved; those prefixes and their
unused-row counts are unchanged.  Their chains are unchanged as labelled
edge sets.  For a column prefix, the number of occupied columns and all
chain lengths are preserved by the increasing row renaming.  The first
failed column prefix is therefore unchanged.  The lexicographic tie
choice is also preserved: a strictly increasing map preserves the
lexicographic order of sorted edge lists, including the transposed lists
used in a column step.  Orientation, first failure and chosen chain all
agree under renaming, and terminality is preserved.

Second, renaming gives a bijection on the branches whose child edges
remain in $[r]\cup S$.  In a row step, the inverse images of the sets
$B$ and $U$ are exactly the corresponding sets in the compressed
rectangle: order comparisons with the old endpoint and occupied-row
membership are preserved.  An image-supported replacement therefore
has a unique inverse injection.  Its column set, cardinality, exclusion
of the original chain, collision status and sign are all preserved.
For a column step the permissible column sets are unchanged; the same
argument, using row membership, applies.  A child cannot encode two
different replacement injections, since its replacement edges are
recovered by removing the fixed context: their columns are disjoint
from the context columns in a row step, and their rows are disjoint
in a column step.  The lower and upper branch families are disjoint
because their cardinalities differ.  Thus no multiplicity is lost.
There may be other branches using rows outside $S$, but these belong to
different history supports and are not part of this correspondence.

Applying the branch bijection successively gives a two-way, sign-preserving
bijection between marked paths with history exactly
$\{r+1,\ldots,r+k\}$ in the compressed rectangle and marked paths with
history exactly $S$ in the large rectangle.  Its inverse simply renames
every state and the initial mark back.  The endpoint $T$ is fixed because
it uses only old rows.  All intermediate states have at most $D$ edges,
so the selector argument applies at every step.

Define $c_k(P,T,j)$ to be the total weight of the marked paths in the
rectangle $[r+k]\times[m]$ whose full history is all $k$ appended rows.
Partitioning the marked paths by their history support, and then by its
cardinality, proves the integer identity
\begin{equation}\label{tx:eq:marked-coordinate}
 \sum_{i=r+1}^{r+L}
       C_{r+L,m}(P\cup\{(i,j)\},T)
    =\sum_{k=1}^{D}\binom{L}{k}c_k(P,T,j).
\end{equation}
The coefficients are independent of $L$.  Its constant term is absent
because each path retains its initial mark in the \emph{history}, not
necessarily in its final state.  This distinction is essential.

\section{Rectangular designs and their extensions}

We now pass from integer reduction coordinates to the fixed field $K$.
The two goals are a design with $\psi(\varnothing)=1$ and an extension
of every degree-$(D-1)$ design; the first will detect $1$, and the second
will make the obstruction survive ordinary-PC multiplication.
Assume now the hypotheses of Lemma~\ref{tx:lem:rectangle-obstruction}, with
$D\geq1$, and put $q=p^a$ and $L=m-r$.  For $1\leq k\leq D<q$ the
integer $\binom{L}{k}$ is zero in $K$.  Indeed, $q\mid L$ and in
characteristic $p$,
\[
 (1+z)^L=(1+z^q)^{L/q};
\]
the right side has no coefficient in these degrees.  Take the square
of size $m$ and an old terminal matching $T\subseteq[r]\times[m]$.
The function
\[
 \theta_T(P)=C_{m,m}(P,T)\quad\text{in }K,
       \qquad P\subseteq[r]\times[m],
\]
is a degree-$D$ rectangular design.  Its row equations are the square
row equations, with the same column set.  Its column equations follow
from the square column equation after deleting the appended-row sum,
which vanishes by~\eqref{tx:eq:marked-coordinate} for $|P|<D$ and an
unused column.  If the column is occupied, every appended insertion
collides, so the omitted sum is already zero.  If $P$ collides, every
insertion still collides and both sides are zero.  This covers all
design equations.

On old terminal inputs, the square reduction is still terminal and
\[
 \theta_T(S)=\delta_{S,T}.
\]
In particular, $\theta_{\varnothing}$ is normalised.  For an arbitrary
$\psi\in E_{D-1}(r,m)$ define explicitly
\[
 \widetilde\psi(P)=
    \sum_{\substack{T\text{ terminal in }[r]\times[m]\\|T|\leq D-1}}
                     \psi(T)\theta_T(P).
\]
This is a member of $E_D(r,m)$ and agrees with $\psi$ on all terminal
inputs of size at most $D-1$.  Apply~\eqref{tx:eq:terminal-expansion} in the
old rectangle to both designs, after restricting the new design to
degree $D-1$.  They agree on every edge set of size at most $D-1$.
Thus the extension exists for every unnormalised design, with no
choice of an unspecified lift.

\section{Why extension is precisely the multiplication rule}
The polynomial spaces $P_{\le0}\subseteq P_{\le1}\subseteq\cdots$
form a \emph{filtration}: nested spaces indexed by the degree allowed.
The corresponding design problem asks for compatible values on each
window. Extending a design means keeping every tested value on edge
sets of size at most $D-1$ and supplying the extra degree-$D$ values.
It is stronger than constructing one normalised design at degree $D$;
the multiplication argument needs an extension for whichever
lower-degree table a given edge produces.

Fix an edge $e$. To test a new line $x_e f$ of degree at most $D$,
start with any degree-$D$ design $\psi$. Force the edge into the design
by defining $s_e\psi(P)=\psi(P\cup\{e\})$. This is a degree-$(D-1)$
design. Extend it back to degree $D$. The extension evaluates every
monomial of $f$ exactly as the original design evaluates that monomial
multiplied by $x_e$, because $f$ has degree at most $D-1$.

The induction hypothesis says that \emph{every} degree-$D$ design kills
earlier lines. It therefore kills $f$ under this chosen extension,
which proves that the original $\psi$ kills $x_e f$. This is the exact
point at which a space of designs succeeds where one isolated witness
would be insufficient.

\section{The invariant through ordinary-PC degree falls}

For a design $\psi$ define a linear functional on the original polynomial
ring by
\[
 \Lambda_\psi\!\left(\sum_\alpha b_\alpha x^\alpha\right)
      =\sum_\alpha b_\alpha\psi(\operatorname{supp}\alpha),
\]
where $\operatorname{supp}\alpha$ is the set of edges with positive
exponent.  This formula is a functional on ordinary polynomials; it
does not replace a PC proof by a proof in a Boolean quotient.
Every degree-$D$ design annihilates each axiom usable in a degree-$D$
proof.  Totalities use~\eqref{tx:eq:matching-design} at $P=\varnothing$;
collisions use vanishing on nonmatchings; and the Boolean equations
use the equality of the supports of $x_e^2$ and $x_e$.

For an edge $e$, put $(s_e\psi)(P)=\psi(P\cup\{e\})$.
If $\psi\in E_D(r,m)$, then $s_e\psi\in E_{D-1}(r,m)$.
Collision vanishing is immediate.  In either totality equation at
$|P|<D-1$, apply the degree-$D$ equation to $P\cup\{e\}$, whose
cardinality is at most $D-1$.  The insertion operations commute,
including repeated insertion of the same edge.  Also, directly from
the monomial formula,
\[
 \Lambda_\psi(x_ef)=\Lambda_{s_e\psi}(f).
\]

Induct through a degree-$D$ ordinary-PC derivation, asserting that every
line is annihilated by \emph{every} member of $E_D(r,m)$.  Input axioms
were handled above, and linear combinations preserve the assertion
even when their highest-degree terms cancel.  Consider multiplication
of an earlier line $f$ by $x_e$.  If $f=0$ there is nothing to prove.
Otherwise ordinary polynomial degree gives
$\deg(x_ef)=\deg f+1$, so $\deg f\leq D-1$.
For a given $\psi\in E_D(r,m)$, extend $s_e\psi$ to a design
$\nu\in E_D(r,m)$ using the preceding construction.  The earlier line
is annihilated by $\nu$ by the inductive assertion.  Since the two
designs agree on all edge sets of size at most $D-1$, and every monomial
of $f$ has support of that size or smaller,
\[
 \Lambda_\psi(x_ef)=\Lambda_{s_e\psi}(f)=\Lambda_\nu(f)=0.
\]
This proves the invariant for multiplication as well.  Finally the
normalised design $\theta_{\varnothing}$ sends the constant $1$ to $1$,
so that constant cannot be the last line of the derivation.

For $D=0$, all displayed axioms have positive degree because both sides
of the rectangle are nonempty.  No nonzero constant axiom is available,
and multiplication of a nonzero polynomial cannot produce a degree-zero
line.  Linear combinations of the available zero lines remain zero.
Thus a degree-zero refutation is impossible directly.  This completes
the ordinary-PC proof in every case.  The NS conclusion follows from
the degree-preserving simulation of a total-product NS certificate by
ordinary PC.
